% Informe del proyecto de generación de código MIPS de Compiscript.
% Compilar con: tectonic -X compile informe.tex
\documentclass[11pt,a4paper]{article}

\usepackage[T1]{fontenc}
\usepackage{lmodern}
\usepackage[spanish,es-tabla,es-noquoting]{babel}
\usepackage[margin=2.5cm]{geometry}
\usepackage{graphicx}
\usepackage{xcolor}
\usepackage{booktabs}
\usepackage{longtable}
\usepackage{array}
\usepackage{tabularx}
\usepackage{listings}
\usepackage{caption}
\usepackage{float}
\usepackage{enumitem}
\usepackage{fancyhdr}
\usepackage{needspace}
\usepackage[hidelinks]{hyperref}
\usepackage{xurl}

% Tolerancia extra para párrafos con nombres largos en \texttt.
\setlength{\emergencystretch}{3em}

\definecolor{fondoCodigo}{RGB}{247,248,250}
\definecolor{bordeCodigo}{RGB}{200,205,214}
\definecolor{comentario}{RGB}{95,120,95}
\definecolor{palabraClave}{RGB}{40,70,150}
\definecolor{cadena}{RGB}{150,60,40}
\definecolor{etiquetaTAC}{RGB}{120,50,120}
\definecolor{verdeUVG}{RGB}{0,105,60}

\lstdefinelanguage{Compiscript}{
  morekeywords={let,const,var,function,class,new,this,return,if,else,while,
    do,for,foreach,in,break,continue,switch,case,default,try,catch,print,
    integer,float,string,boolean,void,null,true,false},
  sensitive=true,
  morecomment=[l]{//},
  morecomment=[s]{/*}{*/},
  morestring=[b]",
}

% Lenguaje intermedio TAC de Compiscript: palabras de control resaltadas.
\lstdefinelanguage{TAC}{
  morekeywords={begin_program,end_program,begin_func,end_func,param_in,param,
    call,call_method,return,declare,goto,if_true,if_false,label,print},
  sensitive=true,
  morecomment=[l]{//},
  morestring=[b]",
}

% Ensamblador MIPS32 para MARS.
\lstdefinelanguage{MIPS}{
  morekeywords={addiu,addu,subu,and,or,sll,li,la,lw,sw,sb,lb,move,jal,jr,j,
    beq,bne,blez,bltz,nop,syscall,mul,div,mfhi,mflo},
  morekeywords=[2]{.text,.data,.globl,.align,.word,.byte,.double},
  keywordstyle=[2]\color{etiquetaTAC}\bfseries,
  sensitive=true,
  morecomment=[l]{\#},
  morestring=[b]",
}

\lstset{
  basicstyle=\ttfamily\footnotesize,
  keywordstyle=\color{palabraClave}\bfseries,
  commentstyle=\color{comentario},
  stringstyle=\color{cadena},
  backgroundcolor=\color{fondoCodigo},
  frame=single,
  rulecolor=\color{bordeCodigo},
  breaklines=true,
  showstringspaces=false,
  extendedchars=true,
  numbers=left,
  numberstyle=\tiny\color{gray},
  numbersep=8pt,
  xleftmargin=18pt,
  framexleftmargin=18pt,
  captionpos=b,
  literate={á}{{\'a}}1 {é}{{\'e}}1 {í}{{\'i}}1 {ó}{{\'o}}1 {ú}{{\'u}}1
           {ñ}{{\~n}}1 {Á}{{\'A}}1 {É}{{\'E}}1 {Ó}{{\'O}}1 {→}{{$\rightarrow$}}1,
}

% Listado sin numeración para los fragmentos pequeños.
\lstdefinestyle{sinnumero}{numbers=none,xleftmargin=4pt,framexleftmargin=4pt}

\pagestyle{fancy}
\fancyhf{}
\fancyhead[L]{\small Fase de Compilación: Generación de Código MIPS}
\fancyhead[R]{\small Compiscript --- ANTLR + C++}
\fancyfoot[C]{\thepage}
\renewcommand{\headrulewidth}{0.4pt}

\captionsetup{font=small,labelfont=bf}
\setlength{\parskip}{0.5em}
\setlength{\parindent}{0pt}

\newcolumntype{L}[1]{>{\raggedright\arraybackslash}p{#1}}

% Encabezado en línea propia para las tablas que siguen (un \paragraph en línea
% quedaría después del flotante).
\newcommand{\subhead}[1]{\par\medskip\noindent\textbf{#1}\par\nopagebreak\smallskip}

\begin{document}

% ------------------------------------------------------------------ CARÁTULA
\begin{titlepage}
\begin{center}

\vspace*{0.2cm}

\includegraphics[width=0.24\textwidth]{img/uvg-logo.png}

\vspace{0.8cm}

{\Large\scshape Universidad del Valle de Guatemala}\\[0.30cm]
{\large Facultad de Ingeniería}\\[0.20cm]
{\large Departamento de Ciencias de la Computación}\\[0.20cm]
{\large Construcción de Compiladores}

\vspace{1.1cm}

\rule{\textwidth}{0.8pt}

\vspace{0.5cm}

{\Huge\bfseries [M3] Fase de Compilación:}\\[0.30cm]
{\Huge\bfseries Generación de Código MIPS}

\vspace{0.5cm}

{\Large Backend MIPS32 y ejecución en MARS para Compiscript}\\[0.30cm]
{\small ANTLR 4.13.2 \ $\bullet$ \ C++17 \ $\bullet$ \ TAC \ $\bullet$ \ MIPS32 \ $\bullet$ \ MARS 4.5 \ $\bullet$ \ IDE web}\\[0.30cm]
{\scriptsize\url{https://github.com/G1LB3T0/Proyecto1_Compis_Generador_Analizador_Lexico/tree/Compiladores2-Proyecto3-MIPS}}

\vspace{0.5cm}

\rule{\textwidth}{0.8pt}

\vspace{1.2cm}

{\large\bfseries Catedrático}\\[0.25cm]
{\large Ing. Pablo Koch}

\vspace{0.9cm}

{\large\bfseries Integrantes}\\[0.3cm]
{\large Luis Gonzalez --- 23353}\\[0.2cm]
{\large Joel Jaquez --- 23369}

\vfill

{\large Guatemala, 8 de noviembre de 2026}\\[0.18cm]
{\small Versión del código documentada: \texttt{a82cd63}}

\end{center}
\end{titlepage}

\tableofcontents
\newpage

% ------------------------------------------------------------ 1 INTRODUCCIÓN
\section{Introducción}

Este informe documenta la fase final del compilador de \textbf{Compiscript}: la
\textbf{generación de código MIPS}. La fase parte del código de tres direcciones
(TAC) y de la tabla de símbolos extendida que produjo la fase anterior, y los
traduce a \textbf{ensamblador MIPS32} que el simulador \textbf{MARS} ensambla a
código máquina y ejecuta. Con ello el compilador recorre, por primera vez, el
camino completo desde un archivo \texttt{.cps} hasta un programa que corre y
produce salida observable.

Como en las fases anteriores, \textbf{todo el compilador está escrito en C++17}.
El backend es una tercera pasada, el \texttt{MipsGenerator}, que se enlaza en la
misma biblioteca nativa que el analizador semántico y el generador de TAC.
Python aparece únicamente como adaptador: el servidor Flask del IDE y un módulo
que invoca a MARS y recoge su salida. Ninguna regla de traducción a MIPS está
implementada fuera de C++.

La fase aporta cuatro piezas, que corresponden a los componentes evaluados:

\begin{itemize}[leftmargin=1.4em]
  \item \textbf{Un generador MIPS32} que traduce cada familia de cuádruplos
        ---datos, expresiones, saltos, funciones, listas, objetos, campos,
        métodos y excepciones--- a instrucciones, junto con un entorno de
        ejecución propio para aritmética mixta, cadenas, impresión y errores.
  \item \textbf{Una convención de llamadas} con registros de activación
        físicos de 64 bytes de cabecera, enlaces de control y de acceso,
        paso de parámetros por pila y retorno en \texttt{\$v0}/\texttt{\$v1}.
  \item \textbf{Un algoritmo de asignación de registros} basado en
        \texttt{getReg}, con descriptores de registros, política LRU y
        \emph{spills} a la pila, con métricas observables.
  \item \textbf{La integración con MARS}: ensamblado, ejecución con límite de
        pasos, volcado del código máquina y una consola de ejecución en el IDE.
\end{itemize}

La validación crece a \textbf{193 casos automatizados}: 115 escenarios C++ de
semántica, TAC y MIPS, 66 programas que se ejecutan de verdad en MARS y cuya
salida se compara con un resultado esperado, y 12 pruebas de integración del
IDE.

% -------------------------------------------------------------- 2 OBJETIVOS
\section{Objetivos}

\subsection{Objetivo general}

Implementar un generador de código objeto para Compiscript que, a partir del
código de tres direcciones y la tabla de símbolos de las fases anteriores,
emita ensamblador MIPS32 correcto y ejecutable en un simulador, manejando
llamadas, retornos, registros de activación y asignación de registros.

\subsection{Objetivos específicos}

\begin{itemize}[leftmargin=1.4em]
  \item Definir una representación física de los valores del lenguaje sobre
        palabras de 32 bits, que conserve el tipo efectivo de cada valor.
  \item Diseñar e implementar una convención de llamadas con prólogo, epílogo,
        enlaces de control y de acceso, paso de parámetros y retorno.
  \item Traducir cada instrucción del TAC a una secuencia MIPS equivalente,
        incluidas funciones anidadas, recursión, clases, herencia, listas y
        \texttt{try/catch}.
  \item Implementar un algoritmo \texttt{getReg} con descriptores de registros
        y \emph{spills}, y exponer sus métricas.
  \item Ejecutar el ensamblador en MARS, distinguir errores de ensamblado y de
        ejecución, y empaquetar el código máquina resultante.
  \item Ampliar el IDE y las baterías de pruebas para compilar, ejecutar y
        verificar los programas generados.
\end{itemize}

% ------------------------------------------------------------ 3 ARQUITECTURA
\section{Arquitectura general}

\subsection{El pipeline}

El flujo de información va del texto fuente al programa ejecutado en MARS. Las
etapas nuevas de esta fase son el \texttt{MipsGenerator} y el adaptador del
simulador:

\begin{center}
\begin{minipage}{0.86\textwidth}
\ttfamily\small
\begin{verbatim}
Código .cps
    |
    v
CompiscriptLexer / CompiscriptParser  (ANTLR, C++)
    |
    v
SemanticAnalyzer  (Visitor C++, primera pasada)
    |-- tipos, ámbitos, tabla de símbolos
    v
TacGenerator      (Visitor C++, segunda pasada)
    |-- cuádruplos tipados + registros de activación
    v
(solo si léxico, sintaxis, semántica y TAC son correctos)
    v
MipsGenerator     (C++, tercera pasada)        <<< NUEVO
    |-- rutinas, frames físicos, getReg, runtime
    v
JSON: diagnósticos + símbolos + TAC + MIPS + frames
    |
    v
Flask (adaptador) -> IDE web (Código MIPS, Consola MIPS)
    |
    v
mips_runtime.py -> MARS 4.5 (ensamblado y ejecución)   <<< NUEVO
\end{verbatim}
\end{minipage}
\end{center}

\subsection{Una tercera pasada sobre cuádruplos}

A diferencia del analizador semántico y del generador de TAC, el
\texttt{MipsGenerator} \textbf{no es un Visitor del árbol}. Recibe el
\texttt{SemanticResult} y el \texttt{TacResult} y trabaja sobre los cuádruplos
estructurados: usa \texttt{op}, \texttt{arg1}, \texttt{arg2} y \texttt{result}
de cada \texttt{TacInstruction}. La representación textual del TAC es solo una
vista para lectura humana; nunca se vuelve a tokenizar.

Para seleccionar instrucciones y aplicar coerciones, el generador necesita el
tipo de cada operando. Por eso esta fase extiende \texttt{TacInstruction} con
tres campos nuevos ---\texttt{result\_type}, \texttt{arg1\_type} y
\texttt{arg2\_type}--- que el generador de TAC llena al emitir cada cuádruplo.

\subsection{Núcleo independiente de la interfaz}

El núcleo sigue exponiendo una sola función pública, \texttt{analyzeSource}, que
devuelve un \texttt{AnalysisResult}. Ahora incluye el bloque \texttt{mips}
(\texttt{MipsResult}) y el indicador \texttt{mips\_ok}. El CLI, las baterías de
pruebas y el IDE usan el mismo servicio, de modo que lo que se prueba es
exactamente lo que se ejecuta.

\subsection{Cuándo se genera código}

El backend \textbf{solo se ejecuta cuando todas las fases anteriores terminaron
sin errores}. Si alguna falla, el resultado MIPS queda marcado como
\texttt{skipped} con una explicación en \texttt{reason}:

\begin{table}[H]
\centering
\small
\begin{tabular}{@{}L{4.3cm}cccc@{}}
\toprule
\textbf{Entrada} & \textbf{Sintaxis} & \textbf{Semántica} & \textbf{TAC} & \textbf{MIPS} \\
\midrule
Error léxico o sintáctico & con errores & omitida & omitido & omitido \\
Error semántico & correcta & con errores & omitido & omitido \\
Error de generación MIPS & correcta & correcta & generado & con errores \\
Programa válido & correcta & correcta & generado & generado \\
\bottomrule
\end{tabular}
\caption{Estado de cada fase según el punto en que falla el programa.}
\end{table}

Un error de generación MIPS ---por ejemplo, un literal entero que no cabe en 32
bits--- tiene código, categoría \emph{generación MIPS}, mensaje, línea y
columna, igual que los diagnósticos de las fases anteriores. Además, los
resultados de \textbf{compilación} y \textbf{ejecución} son distintos: un
programa puede compilar correctamente y fallar al ejecutarse (división por cero,
índice fuera de rango), y el IDE los presenta por separado.

% ------------------------------------------------- 4 TRAZABILIDAD POR COMMITS
\section{Estructura del repositorio y trazabilidad por commits}

El desarrollo de esta fase se realizó en la rama
\texttt{Compiladores2-Proyecto3-MIPS}, a partir del último commit de la fase de
código intermedio (\texttt{4b379ad}). Cada commit corresponde a una unidad
funcional completa y a un solo autor, de modo que la contribución individual
queda registrada en el historial. Los ocho commits se publicaron el 8 de
noviembre de 2026 y suman 26 archivos modificados, 3\,056 líneas agregadas y 108
eliminadas. Las identidades de Git corresponden a los dos integrantes:

\begin{table}[H]
\centering
\small
\begin{tabular}{lll}
\toprule
\textbf{Identidad en Git} & \textbf{Correo} & \textbf{Integrante} \\
\midrule
\texttt{G1LB3T0}  & luis.yo09@hotmail.com & Luis Gonzalez (23353) \\
\texttt{jaq23369} & jaq23369@uvg.edu.gt   & Joel Jaquez (23369) \\
\bottomrule
\end{tabular}
\caption{Identidades de Git utilizadas en la rama.}
\end{table}

\subsection{Historial de la rama}

La tabla siguiente recorre los ocho commits de la fase en orden cronológico y
explica qué archivos introduce cada uno y por qué. Es, en la práctica, el mapa
de lectura del repositorio.

\begin{longtable}{@{}p{1.6cm}L{2.7cm}L{9.1cm}@{}}
\toprule
\textbf{Commit} & \textbf{Autor} & \textbf{Aporte y archivos} \\
\midrule
\endfirsthead
\toprule
\textbf{Commit} & \textbf{Autor} & \textbf{Aporte y archivos} \\
\midrule
\endhead

\texttt{afe1463} & Luis Gonzalez &
\textbf{[ADD] Implementar modelo y generador MIPS de Compiscript.} El corazón
de la fase. Crea \texttt{mips\_generator.h} y \texttt{mips\_generator.cpp}
(1\,012 líneas): extracción de rutinas, \emph{layout} de frames físicos,
\texttt{getReg}, selección de instrucciones y el entorno de ejecución en
ensamblador. En \texttt{model.h} agrega \texttt{MipsFrame}, \texttt{MipsResult},
los tipos de operandos de \texttt{TacInstruction} y \texttt{AnalysisResult.mips}.
3 archivos, +1\,062/$-$0. \\
\addlinespace

\texttt{402a3c2} & Luis Gonzalez &
\textbf{[UPDATE] Preparar semántica y TAC para generación MIPS.} Ajusta las
fases anteriores a lo que el backend necesita: el TAC registra los tipos de
cada operando y conserva el \texttt{this} léxico de funciones anidadas en
métodos; la semántica impide que \texttt{break}/\texttt{continue} atraviesen una
función anidada y exige retorno en \texttt{try} y en \texttt{catch}. Agrega los
escenarios correspondientes a \texttt{compiscript\_tests.cpp}. 4 archivos,
+223/$-$34. \\
\addlinespace

\texttt{50a4dbf} & Luis Gonzalez &
\textbf{[ADD] Integrar backend MIPS en servicio y CLI.} En \texttt{service.cpp}
invoca al \texttt{MipsGenerator} solo si el TAC es válido. En \texttt{json.cpp}
serializa el bloque \texttt{mips} (texto, diagnósticos, métricas y frames). En
\texttt{main.cpp} agrega \texttt{--emit-mips} y \texttt{-o} para guardar el
ensamblador. 4 archivos, +105/$-$15. \\
\addlinespace

\texttt{db2e0fa} & Joel Jaquez &
\textbf{[ADD] Ejecutar y ensamblar MIPS con MARS.} Crea
\texttt{mips\_runtime.py} (invocación de MARS, límite de pasos, volcados y
paquete ZIP), \texttt{scripts/run\_mips.py} (CLI de ejecución) y
\texttt{scripts/setup\_mars.py} (descarga verificada por SHA-256). Integra la
preparación de MARS en los scripts de construcción. 5 archivos, +233/$-$0. \\
\addlinespace

\texttt{03408a6} & Joel Jaquez &
\textbf{[TEST] Validar ejecución MIPS y algoritmos del compilador.} Crea
\texttt{mips\_execution\_tests.py} (372 líneas, 66 casos que se ejecutan en
MARS) y el programa \texttt{algoritmo\_compilador.cps}; registra la batería en
CTest cuando hay Python, Java y MARS. 3 archivos, +464/$-$0. \\
\addlinespace

\texttt{279ddee} & Joel Jaquez &
\textbf{[ADD] Integrar MIPS y consola de ejecución en el IDE.} Agrega a
\texttt{app.py} las rutas \texttt{/api/simulator}, \texttt{/api/run} y
\texttt{/api/build}; al IDE, las pestañas \emph{Código MIPS} y
\emph{Consola MIPS}, el botón \emph{Ejecutar en MIPS} y el estado del
simulador. Crea \texttt{ide\_integration\_tests.py} con 12 pruebas. 5 archivos,
+374/$-$24. \\
\addlinespace

\texttt{bd0d6b6} & Joel Jaquez &
\textbf{[DOC] Documentar arquitectura y generación MIPS.} Crea
\texttt{docs/generacion-mips.md} (399 líneas) con la representación de valores,
la convención de llamadas, la política de \texttt{getReg}, el uso de MARS y las
pruebas; actualiza \texttt{docs/arquitectura-compiscript.md}. 2 archivos,
+465/$-$5. \\
\addlinespace

\texttt{a82cd63} & Luis Gonzalez &
\textbf{[DOC] Actualizar guía de uso y pruebas del proyecto 3.} Reescribe el
\texttt{README.md} con la arquitectura completa, los requisitos (Java y MARS),
la construcción, el uso del CLI, del IDE y de las tres baterías. 1 archivo,
+130/$-$30. \\
\bottomrule
\caption{Trazabilidad de la rama \texttt{Compiladores2-Proyecto3-MIPS}.}
\end{longtable}

\subsection{Lectura del historial}

El orden de los commits refleja el orden natural de construcción de la fase.
Primero se implementa el generador y su modelo de datos (\texttt{afe1463}) y se
ajustan las fases anteriores a lo que necesita (\texttt{402a3c2}); luego se
conecta al servicio y al CLI (\texttt{50a4dbf}) y se ejecuta en el simulador
real (\texttt{db2e0fa}). Después llegan la evidencia de que el código generado
funciona (\texttt{03408a6}), la interfaz (\texttt{279ddee}) y la documentación
(\texttt{bd0d6b6}, \texttt{a82cd63}). La autoría se reparte entre los dos
integrantes: cuatro commits cada uno.

% -------------------------------------------------- 5 ARCHIVOS POR CAPA
\section{Descripción de los archivos por capa}

La fase no modifica la gramática ni el código generado por ANTLR. Las tablas
siguientes listan los archivos nuevos o modificados, con su tamaño actual en
líneas.

\subsection{Núcleo del compilador}

\begin{table}[H]
\centering
\small
\begin{tabular}{@{}L{5.9cm}L{1.9cm}rL{4.5cm}@{}}
\toprule
\textbf{Archivo} & \textbf{Estado} & \textbf{Líneas} & \textbf{Función} \\
\midrule
\texttt{src/mips\_generator.cpp} & nuevo & 1012 &
Tercera pasada: rutinas, frames, \texttt{getReg}, selección de instrucciones y
entorno de ejecución. \\
\texttt{include/.../mips\_generator.h} & nuevo & 19 &
Interfaz pública \texttt{MipsGenerator::\allowbreak generate()}. \\
\texttt{include/.../model.h} & modificado & 217 &
\texttt{MipsFrame}, \texttt{MipsResult} y tipos de operandos del TAC. \\
\texttt{src/tac\_generator.cpp} & modificado & 1314 &
Emite los tipos de operandos y el \texttt{this} léxico. \\
\texttt{src/semantic\_analyzer.cpp} & modificado & 1288 &
Límites de \texttt{break}/\texttt{continue} y retorno en \texttt{try/catch}. \\
\texttt{src/service.cpp} & modificado & 101 &
Encadena semántica, TAC y MIPS. \\
\texttt{src/json.cpp} & modificado & 330 &
Serializa el bloque \texttt{mips}. \\
\texttt{src/main.cpp} & modificado & 96 &
Opciones \texttt{--emit-mips} y \texttt{-o}. \\
\bottomrule
\end{tabular}
\caption{Archivos del núcleo C++ tocados en esta fase.}
\end{table}

\subsection{Simulador, pruebas, IDE y documentación}

\begin{table}[H]
\centering
\small
\begin{tabular}{@{}L{6.5cm}rL{6.1cm}@{}}
\toprule
\textbf{Archivo} & \textbf{Líneas} & \textbf{Función} \\
\midrule
\texttt{mips\_runtime.py} & 125 &
Adaptador de MARS: ejecución, volcados y paquete de código máquina. \\
\texttt{scripts/run\_mips.py} & 66 & CLI para compilar y ejecutar en MARS. \\
\texttt{scripts/setup\_mars.py} & 37 & Descarga y verificación de \texttt{Mars.jar}. \\
\texttt{tests/mips\_execution\_tests.py} & 372 &
66 programas ejecutados en MARS con salida esperada. \\
\texttt{tests/ide\_integration\_tests.py} & 102 & 12 pruebas de las rutas Flask. \\
\texttt{tests/compiscript\_tests.cpp} & 870 & 115 escenarios de semántica, TAC y MIPS. \\
\texttt{examples/.../algoritmo\_\allowbreak compilador.cps} & 79 &
Algoritmos del propio compilador en Compiscript. \\
\texttt{app.py} & 219 & Rutas de ejecución, ensamblado y estado del simulador. \\
\texttt{frontend/static/js/app.js} & 1247 & Pestañas \emph{Código MIPS} y \emph{Consola MIPS}. \\
\texttt{docs/generacion-mips.md} & 399 & Diseño del backend y uso de MARS. \\
\texttt{README.md} & 397 & Compilación, ejecución y pruebas. \\
\bottomrule
\end{tabular}
\caption{Archivos del simulador, pruebas, interfaz y documentación.}
\end{table}

% ------------------------------------------- 6 REPRESENTACIÓN DE VALORES
\section{Representación física de valores}

El destino es \textbf{MIPS32}: las direcciones y las palabras son de 32 bits.
Compiscript tiene tipos estáticos, pero algunas situaciones requieren conocer el
tipo en tiempo de ejecución: un temporal \texttt{\%tN} reciclado puede contener
valores de tipos distintos en instrucciones sucesivas, la variable de un
\texttt{catch} es de tipo desconocido y \texttt{print} debe formatear cada valor
según su tipo. Por eso \textbf{cada valor ocupa un slot de 8 bytes}: una palabra
de carga útil seguida de una palabra de etiqueta de tipo.

\begin{table}[H]
\centering
\small
\begin{tabular}{@{}rlL{9.4cm}@{}}
\toprule
\textbf{Etiqueta} & \textbf{Tipo} & \textbf{Carga útil} \\
\midrule
0 & \texttt{null} & Cero. \\
1 & \texttt{integer} & Entero con signo de 32 bits. \\
2 & \texttt{boolean} & Cero o uno. \\
3 & \texttt{string} & Dirección de una cadena UTF-8 terminada en cero. \\
4 & \texttt{float} & Dirección de un doble de 64 bits alineado a 8 bytes. \\
6 & Lista & Dirección de su cabecera. \\
7 & Objeto & Dirección de su cabecera. \\
8 & Función & Dirección de su código. \\
\bottomrule
\end{tabular}
\caption{Etiquetas de tipo de un slot de valor.}
\end{table}

Los dobles se almacenan \textbf{por referencia} para que el slot mantenga el
mismo tamaño que los demás valores. Los literales flotantes residen en
\texttt{.data}; los resultados calculados se alojan en memoria dinámica. La
ampliación \texttt{integer} $\rightarrow$ \texttt{float} se realiza al almacenar en
un destino flotante, al pasar parámetros y al devolver valores de ese tipo.

\subhead{Rangos y aritmética.}

\begin{itemize}[leftmargin=1.4em]
  \item Los literales enteros deben estar entre $-2\,147\,483\,648$ y
        $2\,147\,483\,647$; si no, el backend emite \texttt{MIPS\_INTEGER\_RANGE}.
  \item La suma, resta, multiplicación y negación enteras conservan los 32 bits
        bajos del resultado.
  \item Los dobles deben ser finitos: un literal fuera de rango se diagnostica al
        compilar (\texttt{MIPS\_FLOAT\_RANGE}) y un resultado no finito se
        reporta al ejecutar, en lugar de propagarse como infinito o NaN.
  \item La división produce un resultado flotante; la división o el resto por
        cero producen un error de ejecución capturable con \texttt{try/catch}.
\end{itemize}

\subhead{Cadenas, listas y objetos.}

\begin{itemize}[leftmargin=1.4em]
  \item Las cadenas son bytes UTF-8. Las secuencias \texttt{\textbackslash n},
        \texttt{\textbackslash r}, \texttt{\textbackslash t}, \texttt{\textbackslash 0},
        \texttt{\textbackslash\textquotedbl} y \texttt{\textbackslash\textbackslash} se
        decodifican al construir el literal. La salida usa el servicio 15 de MARS
        (escritura en descriptor) para no partir los caracteres multibyte.
  \item Una lista tiene una cabecera de 8 bytes con la longitud en el offset 0;
        el elemento $i$ está en \texttt{base + 8 + 8$\cdot$i}. Una lista anidada
        guarda referencias a otras listas.
  \item Un objeto tiene una cabecera de 8 bytes con una referencia al nombre de
        la clase; un campo está en \texttt{base + 8 + offset\_TAC}, y los campos de
        la clase base preceden a los de la derivada. Los inicializadores de
        campos se ejecutan para cada instancia antes del constructor.
\end{itemize}

La memoria dinámica se obtiene con el servicio 9 (\texttt{sbrk}) del simulador,
se inicializa en cero y se conserva durante toda la ejecución; esta fase no
incorpora un recolector de basura.

% ---------------------------------------- 7 CONVENCIÓN DE LLAMADAS
\section{Convención de llamadas y registros de activación}

\subsection{Organización del programa en rutinas}

Antes de emitir \texttt{.text}, el generador recorre los cuádruplos y separa
\textbf{rutinas}: el programa principal (\texttt{main}), cada función o método
delimitado por \texttt{begin\_func}/\texttt{end\_func} y un inicializador
\texttt{\_\_init\_<Clase>} por cada clase. Una pila de rutinas abiertas permite
que una función declarada dentro de otra se extraiga como rutina propia, sin
que su cuerpo se ejecute en el punto de la declaración. Así las funciones pueden
declararse antes o después de sus llamadas.

\subsection{Registro de activación físico}

Cada rutina reserva un frame alineado a 8 bytes. \texttt{\$fp} apunta a su
inicio y \texttt{\$sp} delimita el espacio de la activación y de las llamadas
que realiza. La cabecera ocupa 64 bytes:

\begin{table}[H]
\centering
\small
\begin{tabular}{@{}rL{10.4cm}@{}}
\toprule
\textbf{Offset desde \texttt{\$fp}} & \textbf{Contenido} \\
\midrule
0 & \texttt{\$fp} anterior: \textbf{enlace de control}. \\
4 & Frame del padre léxico: \textbf{enlace de acceso}. \\
8 & Dirección de retorno \texttt{\$ra}. \\
12--40 & Valores preservados de \texttt{\$s0} a \texttt{\$s7}. \\
44 & Manejador de errores activo al entrar en la rutina. \\
48 & Carga útil del receptor \texttt{this}, cuando existe. \\
52 & Etiqueta del receptor. \\
56--63 & Reservado. \\
64 en adelante & Slots de 8 bytes para parámetros copiados, locales y temporales. \\
\bottomrule
\end{tabular}
\caption{Cabecera del registro de activación físico.}
\end{table}

El tamaño del frame es $64 + 8 \cdot n$, donde $n$ es el número de nombres
distintos (variables locales, parámetros y temporales) que la rutina usa. Cada
nombre recibe un slot estable en el orden en que aparece, de modo que un valor
sacado de un registro siempre tiene dónde guardarse. El generador publica estos
datos en \texttt{MipsFrame}: \texttt{frame\_size}, \texttt{parameter\_bytes},
\texttt{local\_bytes}, \texttt{temporary\_bytes} y \texttt{saved\_register\_bytes}
(los 64 bytes de cabecera).

\subsection{Responsabilidades del llamante}

\begin{enumerate}[leftmargin=1.6em]
  \item Vacía los registros modificados a sus slots (\texttt{flush}).
  \item Reserva $8 \cdot n$ bytes en la pila y escribe allí los argumentos, de
        izquierda a derecha; en un método, el receptor precede a los argumentos
        explícitos. Si el parámetro formal es \texttt{float}, aplica la
        promoción antes de escribirlo.
  \item Coloca en \texttt{\$a0} el frame del \textbf{padre léxico} de la rutina
        destino, recorriendo la cadena de enlaces de acceso desde su propio
        frame.
  \item Salta con \texttt{jal} y, al volver, retira el área de argumentos y
        copia \texttt{\$v0}/\texttt{\$v1} al destino del cuádruplo.
\end{enumerate}

\subsection{Responsabilidades del llamado}

El prólogo y el epílogo que emite \texttt{Generator::routine} son, para una
rutina cuyo frame mide \texttt{F} bytes:

\begin{lstlisting}[language=MIPS,caption={Prólogo y epílogo generados para cada rutina.}]
fn_etiqueta:
  addiu $sp, $sp, -F       # reserva el frame
  sw $fp, 0($sp)           # enlace de control
  sw $a0, 4($sp)           # enlace de acceso recibido del llamante
  sw $ra, 8($sp)           # dirección de retorno
  sw $s0, 12($sp)          # ... $s0-$s7 en 12..40
  lw $t8, __handler
  sw $t8, 44($sp)          # manejador de errores de entrada
  move $fp, $sp
  sw $zero, 64($fp)        # cada slot inicia en null
  ...
  # cuerpo: TAC traducido
  ...
fn_etiqueta__return:
  lw $t8, 44($fp)
  sw $t8, __handler        # restaura el manejador de entrada
  move $sp, $fp
  lw $ra, 8($sp)
  lw $s0, 12($sp)          # ... restaura $s0-$s7
  lw $fp, 0($sp)
  addiu $sp, $sp, F
  jr $ra
  nop
\end{lstlisting}

Los parámetros se leen desde el área del llamante, que está justo encima del
frame: el parámetro $i$ está en \texttt{\$fp + F + 8$\cdot$i} (el índice incluye
el receptor implícito en los métodos). Los métodos e inicializadores copian el
receptor a \texttt{\$s0} y al offset 48. El retorno utiliza \texttt{\$v0} para la
carga útil y \texttt{\$v1} para la etiqueta; todo \texttt{return} salta a la
etiqueta \texttt{\_\_return} común, y una rutina que termina sin
\texttt{return} devuelve \texttt{null}.

\subsection{Acceso a datos no locales}

Cuando una instrucción usa una variable de una función exterior, el generador
calcula su dirección con la rutina \texttt{ancestor}: parte de \texttt{\$fp} y
sigue el enlace de acceso (offset 4) tantas veces como niveles léxicos haya
entre la rutina actual y la propietaria del símbolo; luego suma el offset del
slot en el frame del propietario.

\begin{lstlisting}[language=MIPS,style=sinnumero,caption={Acceso a una variable dos niveles léxicos afuera.}]
  move $t8, $fp
  lw $t8, 4($t8)           # frame del padre léxico
  lw $t8, 4($t8)           # frame del abuelo léxico
  addiu $t8, $t8, 72       # slot de la variable capturada
  lw $s1, 0($t8)           # carga útil
  lw $s2, 4($t8)           # etiqueta
\end{lstlisting}

El \textbf{padre de definición puede ser distinto del llamante}: una llamada
entre funciones hermanas o una llamada recursiva recibe en \texttt{\$a0} el
entorno exterior que le corresponde léxicamente, no el frame de quien la
invoca. Una función declarada dentro de un método accede a su \texttt{this}
léxico copiándolo del frame del método, sin convertirse en un método con un
argumento receptor adicional.

\subsection{Recursividad}

Cada llamada crea una activación independiente: los temporales y variables de
una llamada recursiva viven en su propio frame y no comparten almacenamiento con
la activación que la invocó. Los registros \texttt{\$s0}--\texttt{\$s7} y la
dirección de retorno se preservan en la cabecera, de modo que los valores que
deben sobrevivir a una llamada (por ejemplo, \texttt{n} en
\texttt{n * factorial(n - 1)}) se recuperan intactos al volver.

% ------------------------------------------ 8 SELECCIÓN DE INSTRUCCIONES
\section{Selección de instrucciones}

Cada cuádruplo se traduce en \texttt{Generator::instruction}. Antes de cada
traducción se emite un comentario \texttt{\# TAC n: ...} con el texto del
cuádruplo original, de modo que el ensamblador puede leerse en paralelo con el
TAC.

\begin{longtable}{@{}L{3.6cm}L{10.6cm}@{}}
\toprule
\textbf{Operación TAC} & \textbf{Traducción MIPS} \\
\midrule
\endfirsthead
\toprule
\textbf{Operación TAC} & \textbf{Traducción MIPS} \\
\midrule
\endhead
\texttt{declare} & Carga el valor por defecto del tipo (\texttt{0},
\texttt{false} o \texttt{null}) y lo escribe en el destino. \\
\texttt{assign} & Copia carga útil y etiqueta del origen al destino, con
promoción a \texttt{float} si el destino lo requiere. \\
\texttt{add}, \texttt{sub}, \texttt{mul}, \texttt{div}, \texttt{mod} &
Operandos en \texttt{\$a0}--\texttt{\$a3}; \texttt{jal \_\_add} (etc.). La rutina
auxiliar decide aritmética entera, flotante o concatenación según las
etiquetas. \\
\texttt{eq}, \texttt{ne}, \texttt{lt}, \texttt{le}, \texttt{gt}, \texttt{ge} &
\texttt{jal \_\_eq} (etc.), sobre \texttt{\_\_compare}: compara números mixtos y
cadenas por contenido y orden lexicográfico. \\
\texttt{neg}, \texttt{not} & \texttt{jal \_\_neg} / \texttt{\_\_not}. \\
\texttt{label} & Vacía registros y emite la etiqueta. \\
\texttt{goto} & Vacía registros y emite \texttt{j L}; si el salto sale de un
\texttt{try}, retira antes los manejadores. \\
\texttt{if\_true}, \texttt{if\_false} & Carga la condición y emite
\texttt{bne}/\texttt{beq \$a0, \$zero, L}. \\
\texttt{print} & \texttt{jal \_\_print} y salto de línea con el servicio 4. \\
\texttt{param}, \texttt{call}, \texttt{call\_method} & Protocolo de llamada de
la sección~7. \\
\texttt{param\_in} & Lee el argumento desde \texttt{\$fp + F + 8$\cdot$i}. \\
\texttt{return} & Valor en \texttt{\$v0}/\texttt{\$v1}, promoción al tipo de
retorno y salto al epílogo. \\
\texttt{new\_list}, \texttt{length} & \texttt{jal \_\_new\_list}; lectura de la
longitud en la cabecera con verificación de \texttt{null}. \\
\texttt{load\_index}, \texttt{store\_index} & \texttt{jal \_\_index}, que
verifica límites y devuelve la dirección del elemento. \\
\texttt{new\_object} & Reserva la cabecera y los campos con
\texttt{\_\_alloc\_zero}, guarda el nombre de la clase y llama al inicializador
de la clase. \\
\texttt{load\_field}, \texttt{store\_field} & Verifica \texttt{null} y accede a
\texttt{base + 8 + offset}. \\
\texttt{try\_begin}, \texttt{try\_end}, \texttt{catch} & Apila o retira un
registro de manejador (sección~10). \\
\bottomrule
\caption{Traducción de cada familia de cuádruplos.}
\end{longtable}

Las operaciones \texttt{begin\_program}, \texttt{end\_program},
\texttt{capture} y \texttt{static\_link} no emiten código: su información ya
quedó reflejada en la organización de rutinas y en los enlaces de acceso. Una
operación sin traducción produce el diagnóstico
\texttt{MIPS\_UNSUPPORTED\_TAC} en lugar de un ensamblador incompleto.

\subsection{Ejemplo: una suma}

Para el cuádruplo \texttt{\%t0 = \%a\#2 add \%b\#3} dentro de una función, con
los registros vacíos, el generador emite:

\begin{lstlisting}[language=MIPS,caption={Traducción de una suma (offsets ilustrativos).}]
  # TAC 4: %t0 = %a#2 + %b#3
  addiu $t8, $fp, 64       # getReg(%a#2): par $s1/$s2
  lw $s1, 0($t8)
  lw $s2, 4($t8)
  move $a0, $s1
  move $a1, $s2
  addiu $t8, $fp, 72       # getReg(%b#3): par $s3/$s4
  lw $s3, 0($t8)
  lw $s4, 4($t8)
  move $a2, $s3
  move $a3, $s4
  jal __add                # aritmética según etiquetas
  nop
  move $s1, $v0            # getReg(%t0): el destino queda en registro
  move $s2, $v1            # y se marca como modificado
\end{lstlisting}

El resultado no se escribe en memoria de inmediato: queda en el par
\texttt{\$s1}/\texttt{\$s2} marcado como modificado y solo se vuelca a su slot
si el par se necesita para otro valor o antes de una llamada o un salto.

\subsection{Entorno de ejecución}

Al final de \texttt{.text} el generador agrega un conjunto de rutinas auxiliares
escritas en ensamblador. Concentran las operaciones que dependen de la
etiqueta de tipo y mantienen compacto el código de cada cuádruplo:

\begin{table}[H]
\centering
\small
\begin{tabular}{@{}L{4.4cm}L{9.8cm}@{}}
\toprule
\textbf{Rutina} & \textbf{Responsabilidad} \\
\midrule
\texttt{\_\_alloc\_zero} & Reserva memoria con \texttt{sbrk}, alineada a 8 y en cero. \\
\texttt{\_\_new\_list}, \texttt{\_\_index} & Creación de listas y cálculo de
dirección con verificación de límites. \\
\texttt{\_\_add} \dots{} \texttt{\_\_mod}, \texttt{\_\_to\_float} &
Aritmética entera y flotante, promoción y empaquetado de dobles. \\
\texttt{\_\_compare}, \texttt{\_\_eq} \dots{} \texttt{\_\_ge} & Comparación
numérica mixta y de cadenas. \\
\texttt{\_\_concat} & Concatenación de cadenas UTF-8. \\
\texttt{\_\_print}, \texttt{\_\_print\_list} & Impresión de cada tipo, incluidas
listas anidadas, objetos y \texttt{null}. \\
\texttt{\_\_raise}, \texttt{\_\_uncaught} & Lanzamiento de errores y
terminación con código 2 si nadie los captura. \\
\bottomrule
\end{tabular}
\caption{Rutinas auxiliares del entorno de ejecución.}
\end{table}

El segmento \texttt{.data} contiene el manejador activo, un slot por variable
global (\texttt{\_\_global\_<id>}), los literales de cadena
(\texttt{\_\_str\_n}) y flotantes (\texttt{\_\_flt\_n}) sin duplicados, y los
mensajes del entorno.

% ------------------------------------------ 9 ASIGNACIÓN DE REGISTROS
\section{Asignación de registros}

\subsection{Registros disponibles}

\begin{table}[H]
\centering
\small
\begin{tabular}{@{}L{3.8cm}L{10.4cm}@{}}
\toprule
\textbf{Registros} & \textbf{Uso} \\
\midrule
\texttt{\$s1/\$s2}, \texttt{\$s3/\$s4}, \texttt{\$s5/\$s6} & Tres pares de la
caché de valores: carga útil y etiqueta. \\
\texttt{\$s0} & Receptor \texttt{this} de la rutina. \\
\texttt{\$a0}--\texttt{\$a3} & Operandos de las rutinas auxiliares; \texttt{\$a0}
transporta además el enlace de acceso en las llamadas. \\
\texttt{\$v0}/\texttt{\$v1} & Resultado de rutinas y funciones. \\
\texttt{\$t0}--\texttt{\$t9} & Registros de trabajo para direcciones y rutinas
auxiliares. \\
\texttt{\$fp}, \texttt{\$sp}, \texttt{\$ra} & Frame, pila y retorno. \\
\bottomrule
\end{tabular}
\caption{Asignación fija de los registros MIPS.}
\end{table}

\subsection{Política de \texttt{getReg}}

El generador mantiene un \textbf{descriptor por par de registros}: qué valor
contiene, si está modificado respecto de su slot (\emph{dirty}) y cuándo se usó
por última vez. \texttt{getReg(v)} aplica, en orden:

\begin{enumerate}[leftmargin=1.6em]
  \item Si algún par ya contiene \texttt{v}, lo reutiliza sin emitir carga.
  \item Si hay un par vacío, lo elige.
  \item Si no, elige el \textbf{menos recientemente usado} (LRU).
  \item Si el par elegido está modificado, emite un \emph{spill}: guarda carga
        útil y etiqueta en el slot estable del valor anterior.
  \item Si \texttt{v} se va a leer, lo carga desde su slot o desde el literal.
\end{enumerate}

\begin{lstlisting}[language=C++,caption={Núcleo de \texttt{getReg} en \texttt{mips\_generator.cpp}.}]
int Generator::getReg(const std::string &v, bool reading, int avoid) {
  for (int n = 0; n < 3; ++n)
    if (cache[n].value == v && !v.empty()) { cache[n].age = ++age; ... return n; }
  int n = -1;
  for (int j = 0; j < 3; ++j) if (j != avoid && cache[j].value.empty()) { n = j; break; }
  if (n < 0) for (int j = 0; j < 3; ++j)
    if (j != avoid && (n < 0 || cache[j].age < cache[n].age)) n = j;
  spill(n); cache[n] = {v, false, ++age};
  if (reading) load(v, p(n), k(n));
  return n;
}
\end{lstlisting}

El parámetro \texttt{avoid} impide que el segundo operando de una operación
binaria desaloje al primero. Antes de cada llamada, etiqueta y salto se ejecuta
\texttt{flush}: se vuelcan los pares modificados y se invalidan los
descriptores, de modo que la memoria refleja los valores que deben estar
disponibles en cualquier camino de control que converja en ese punto.

\subsection{Relación con los temporales del TAC}

La asignación física es independiente del reciclaje de nombres \texttt{\%tN} que
realiza el TAC. Un mismo nombre temporal puede representar valores distintos en
instrucciones sucesivas; como la etiqueta de tipo viaja con el valor, el
reciclaje no confunde al backend. Cada nombre tiene un único slot por frame, así
que la presión de registros se resuelve con \emph{spills} y no requiere
reservar memoria adicional.

\subsection{Métricas}

\texttt{MipsResult} expone \texttt{instruction\_count} (instrucciones del
ensamblador generado, incluido el entorno), \texttt{register\_allocations}
(selecciones de pares hechas por \texttt{getReg}) y \texttt{spills} (guardados
efectivos). Estas métricas describen el código estático; no son el número
dinámico de instrucciones que ejecuta MARS. El IDE las muestra en la pestaña
\emph{Código MIPS}.

Esta es una asignación \textbf{local y conservadora con respaldo en pila}: no
intenta una asignación global por coloreo de grafos, que el enunciado no exige,
y privilegia la corrección en llamadas y en los puntos donde convergen caminos
de control.

% ------------------------------------------ 10 ERRORES Y EXCEPCIONES
\section{Errores de ejecución y \texttt{try/catch}}

La fase anterior conservaba \texttt{try/catch} como operaciones abstractas
porque no existía un ABI de excepciones. Esta fase lo define. La palabra
\texttt{\_\_handler} en \texttt{.data} apunta al manejador activo. Un
\texttt{try\_begin} reserva 64 bytes en la pila y guarda allí:

\begin{itemize}[leftmargin=1.4em]
  \item el manejador anterior (para encadenar manejadores anidados);
  \item la etiqueta del \texttt{catch};
  \item \texttt{\$fp} y \texttt{\$sp} en el momento de entrar al \texttt{try};
  \item los registros \texttt{\$s0}--\texttt{\$s7}.
\end{itemize}

Cuando una rutina auxiliar detecta un error (división por cero, índice fuera de
rango, referencia nula, resultado no finito) salta a \texttt{\_\_raise}, que
restaura \texttt{\$fp}, \texttt{\$sp} y los registros del manejador activo, deja
el mensaje como valor del \texttt{catch} y salta a su etiqueta, aunque el error
haya ocurrido varias llamadas más abajo. Si no hay manejador,
\texttt{\_\_uncaught} imprime \emph{Error de ejecución:} seguido del mensaje y
termina con el servicio 17 y código 2.

Para que los manejadores no queden colgados, cada frame guarda el manejador de
entrada (offset 44) y lo restaura en el epílogo, y los saltos que salen de un
\texttt{try} ---\texttt{return}, \texttt{break}, \texttt{continue}--- retiran
antes los registros de manejador que abandonan.

% ------------------------------------------ 11 MARS
\section{Integración con MARS y código máquina}

\subsection{Preparación del simulador}

\texttt{scripts/setup\_mars.py} descarga el JAR oficial de MARS 4.5 a
\texttt{tools/Mars.jar} y verifica su SHA-256. Las variables \texttt{MARS\_JAR}
y \texttt{JAVA} permiten usar una instalación existente. Los scripts de
construcción preparan MARS antes de configurar CMake.

\subsection{Ejecución}

\texttt{mips\_runtime.py} ejecuta MARS sin interfaz gráfica:

\begin{lstlisting}[language=bash,caption={Comando de referencia de MARS.}]
java -jar tools/Mars.jar nc sm ae65 se70 1000000 programa.asm
\end{lstlisting}

\texttt{nc} elimina el aviso inicial, \texttt{sm} comienza en la etiqueta global
\texttt{main}, \texttt{ae65} y \texttt{se70} asignan códigos de salida distintos
a los errores de ensamblado y de simulación, y el entero limita los pasos
ejecutados (entre 32 y 10\,000\,000), lo que protege al IDE frente a ciclos
infinitos. El reporte distingue \texttt{assembly} de \texttt{execution} y
conserva \texttt{stdout}, \texttt{stderr}, \texttt{returncode},
\texttt{timed\_out} y \texttt{step\_limit}.

\subsection{Código máquina}

El adaptador puede ensamblar sin ejecutar y empaquetar el resultado en un ZIP:

\begin{table}[H]
\centering
\small
\begin{tabular}{@{}L{4.0cm}L{10.2cm}@{}}
\toprule
\textbf{Archivo} & \textbf{Contenido} \\
\midrule
\texttt{program.asm} & Ensamblador del programa y de su entorno de ejecución. \\
\texttt{program.text.bin} & Volcado binario del segmento de instrucciones
ensambladas. \\
\texttt{program.data.bin} & Volcado binario del segmento de datos iniciales. \\
\texttt{manifest.json} & Arquitectura, simulador, direcciones y tamaños de
segmentos, comando de ejecución. \\
\texttt{LEEME.txt} & Instrucciones para abrir y ejecutar el programa en MARS. \\
\bottomrule
\end{tabular}
\caption{Contenido del paquete de código máquina.}
\end{table}

Los binarios son volcados sin cabecera ELF; para reproducir la ejecución se usa
\texttt{program.asm}, que permite a MARS reconstruir direcciones, datos y
código de forma coherente.

% ------------------------------------------ 12 IDE
\section{Interfaz gráfica (IDE)}

El IDE es la misma aplicación web servida por Flask de las fases anteriores. El
botón principal ahora se llama \emph{Compilar programa} y la barra de fases
muestra \emph{ANTLR $\rightarrow$ semántica $\rightarrow$ TAC $\rightarrow$
ensamblador MIPS32 $\rightarrow$ MARS}. La fase agrega:

\begin{itemize}[leftmargin=1.4em]
  \item \textbf{Pestaña \emph{Código MIPS}}: el ensamblador generado, sus
        métricas (instrucciones, selecciones de registros y \emph{spills}), la
        tabla de frames físicos y botones para copiar y descargar el
        \texttt{.asm} o el ZIP de código máquina.
  \item \textbf{Botón \emph{Ejecutar en MIPS} y pestaña \emph{Consola MIPS}}:
        ensambla y ejecuta el programa en MARS y muestra la salida estándar, los
        errores de ensamblado o de ejecución y el código de retorno, separados
        de los diagnósticos del compilador.
  \item \textbf{Estado del simulador}: indica si Java y \texttt{Mars.jar} están
        disponibles antes de intentar ejecutar.
\end{itemize}

\begin{table}[H]
\centering
\small
\begin{tabular}{@{}L{3.6cm}L{10.6cm}@{}}
\toprule
\textbf{Ruta} & \textbf{Función} \\
\midrule
\texttt{POST /api/analyze} & Compila y devuelve el JSON de todas las fases. \\
\texttt{GET /api/simulator} & Disponibilidad de Java y MARS. \\
\texttt{POST /api/run} & Compila y ejecuta en MARS; devuelve salida y errores. \\
\texttt{POST /api/build} & Ensambla y devuelve el ZIP de código máquina. \\
\texttt{POST /api/tests} & Ejecuta la batería C++ desde la interfaz. \\
\bottomrule
\end{tabular}
\caption{Rutas del adaptador Flask.}
\end{table}

Si el programa tiene errores léxicos, sintácticos o semánticos, la tarjeta
\emph{Código MIPS} indica \emph{Omitido} y la ejecución no se intenta: el IDE
nunca envía a MARS código de un programa rechazado.

% ----------------------------------------------------- 13 BATERÍA DE PRUEBAS
\section{Batería de pruebas}

La fase se verifica en tres niveles complementarios, con \textbf{193 casos
automatizados}:

\begin{table}[H]
\centering
\small
\begin{tabular}{@{}L{5.0cm}rL{7.6cm}@{}}
\toprule
\textbf{Batería} & \textbf{Casos} & \textbf{Qué verifica} \\
\midrule
\texttt{compiscript\_tests.cpp} & 115 & Semántica, TAC y omisión de MIPS ante
errores, con el código de diagnóstico exacto. \\
\texttt{mips\_execution\_tests.py} & 66 & Genera ensamblador desde
\texttt{.cps}, lo ejecuta en MARS y compara la salida con un resultado
esperado: 54 válidos y 12 errores esperados. \\
\texttt{ide\_integration\_tests.py} & 12 & Rutas Flask con el compilador y MARS
reales. \\
\midrule
\textbf{Total} & \textbf{193} & \\
\bottomrule
\end{tabular}
\caption{Baterías de pruebas del proyecto.}
\end{table}

\subsection{Pruebas de ejecución en MARS}

La batería \texttt{mips\_execution\_tests.py} \textbf{no interpreta Compiscript
ni TAC}: los resultados observados proceden del ensamblador ejecutado por MARS.
Las salidas esperadas se calcularon a partir de los programas fuente, no del
TAC, y cada caso válido exige además un volcado binario no vacío de
\texttt{.text}, para comprobar que el ensamblado produjo código máquina. Los
casos negativos distinguen un diagnóstico esperado de un fallo inesperado del
ensamblador.

\begin{table}[H]
\centering
\small
\begin{tabular}{@{}L{4.6cm}rL{8.0cm}@{}}
\toprule
\textbf{Categoría} & \textbf{Casos} & \textbf{Ejemplos} \\
\midrule
Proceso completo & 1 & Programa vacío. \\
Expresiones & 5 & Precedencia, promoción flotante, comparaciones mixtas, resto
flotante. \\
Control de flujo & 6 & Cortocircuito con efectos laterales, ciclos,
\texttt{break}/\texttt{continue} anidados, \texttt{switch}. \\
Llamadas y retornos & 8 & Factorial, Fibonacci, recursión mutua, seis
parámetros, retorno flotante. \\
Asignación de registros & 2 & Temporales vivos alrededor de llamadas, presión
con doce argumentos. \\
Orden de evaluación & 3 & Argumentos de izquierda a derecha, operando
conservado ante efecto lateral. \\
Ámbitos y closures & 7 & Sombreado, dos niveles léxicos, funciones hermanas,
\texttt{this} en función anidada. \\
Cadenas & 4 & Concatenación, orden lexicográfico, escapes, UTF-8. \\
Listas & 5 & \texttt{foreach}, listas anidadas, paso por referencia. \\
Clases y objetos & 7 & Constructores, herencia, inicializadores, receptor
preservado. \\
Algoritmos del compilador & 1 & \texttt{algoritmo\_compilador.cps}. \\
Errores de compilación & 7 & Léxico, sintáctico y semántico impiden generar
ensamblador; aridad. \\
Errores de generación & 1 & Literal entero fuera de rango MIPS32. \\
Errores de ejecución & 4 & División por cero, índices fuera de rango, referencia
nula. \\
Excepciones & 5 & \texttt{try/catch} entre llamadas, manejadores anidados,
\texttt{return} y \texttt{continue} dentro de \texttt{try}. \\
\midrule
\textbf{Total} & \textbf{66} & \\
\bottomrule
\end{tabular}
\caption{Casos de ejecución en MARS por categoría.}
\end{table}

\subsection{Prueba de fuego: algoritmos del propio compilador}

\texttt{algoritmo\_compilador.cps} adapta a Compiscript dos algoritmos del
repositorio y recorre con ellos el pipeline completo hasta MARS:

\begin{itemize}[leftmargin=1.4em]
  \item \texttt{DFABuilder::bestRule} (\texttt{backend/src/dfa\_builder.cpp}):
        elige la regla aceptante de menor prioridad y usa cortocircuito para
        evitar una lectura fuera de rango.
  \item \texttt{TacGenerator::newTemp} y \texttt{releaseTemp}
        (\texttt{tac\_generator.cpp}): administran temporales vivos con una pila
        de libres y contadores de creación, reutilización y máximo vivo.
\end{itemize}

\begin{lstlisting}[language=Compiscript,caption={Fragmento de \texttt{algoritmo\_compilador.cps}.}]
function mejorRegla(posiciones: integer[], reglas: integer[], total: integer): integer {
  let mejor: integer = -1;
  foreach (posicion in posiciones) {
    if (posicion < total && reglas[posicion] >= 0) {
      if (mejor == -1 || reglas[posicion] < mejor) {
        mejor = reglas[posicion];
      }
    }
  }
  return mejor;
}
\end{lstlisting}

La adaptación usa listas y una clase; la pila tiene capacidad fija de ocho
posiciones porque el lenguaje no define \texttt{push}, \texttt{pop} ni listas
redimensionables. No es autocompilación de los archivos C++: la sintaxis
Compiscript sigue siendo el contrato de entrada. La salida esperada,
\texttt{2, -1, 7, 0, 1, 0, 2, 2, 3, 2, 3} (un valor por línea), se compara con la
de MARS.

\subsection{Ejecución de las baterías}

\begin{lstlisting}[language=bash,caption={Ejecución de las tres baterías.}]
ctest --test-dir build/compiscript --output-on-failure
python3 backend/tests/mips_execution_tests.py
python3 backend/tests/mips_execution_tests.py --filter "Llamadas" --json
python3 backend/tests/ide_integration_tests.py
\end{lstlisting}

CTest registra la batería MIPS cuando al configurar CMake encuentra Python, Java
y \texttt{tools/Mars.jar}. Las pruebas del IDE requieren instalar antes
\texttt{requirements.txt}.

% ------------------------------------------------ 14 COMPILACIÓN Y EJECUCIÓN
\section{Compilación y ejecución}

\subsection{Requisitos}

\begin{itemize}[leftmargin=1.4em]
  \item CMake 3.20 o superior.
  \item Compilador compatible con C++17.
  \item Java, para la herramienta generadora de ANTLR y para el simulador MARS.
  \item Python 3.10 o superior con Flask, para el IDE y el adaptador de MARS.
\end{itemize}

\subsection{Construcción}

Un solo script descarga el JAR de ANTLR si hace falta, regenera el lexer, el
parser y el Visitor, prepara y verifica MARS, compila el \emph{runtime} C++,
construye \texttt{compiscript\_cli} y \texttt{compiscript\_tests} y ejecuta las
baterías C++ y de ejecución MIPS:

\begin{lstlisting}[language=bash,caption={macOS o Linux.}]
python3 scripts/setup_mars.py
./scripts/build_compiscript.sh
\end{lstlisting}

\begin{lstlisting}[language=bash,caption={Windows PowerShell.}]
py scripts/setup_mars.py
.\scripts\build_compiscript.ps1
\end{lstlisting}

\subsection{Ejecución del IDE}

\begin{lstlisting}[language=bash,caption={Servidor del IDE.}]
python app.py
\end{lstlisting}

El IDE queda disponible en \texttt{http://127.0.0.1:5050}. Se abre un ejemplo,
se pulsa \emph{Compilar programa}, se revisa la pestaña \emph{Código MIPS} y se
pulsa \emph{Ejecutar en MIPS} para ver la salida en la \emph{Consola MIPS}.

\subsection{Uso directo del CLI}

\begin{lstlisting}[language=bash,caption={Generación y ejecución desde la línea de comandos.}]
./compiscript_cli --file programa.cps --emit-mips -o programa.asm
python3 scripts/run_mips.py programa.cps -o programa.asm
python3 scripts/run_mips.py programa.asm
python3 scripts/run_mips.py programa.cps --assemble-only --package programa.zip
\end{lstlisting}

Sin \texttt{--emit-mips}, el CLI conserva la salida JSON con los bloques
\texttt{tokens}, \texttt{tree}, \texttt{semantic}, \texttt{tac} y \texttt{mips}
(\texttt{text}, \texttt{diagnostics}, \texttt{frames}, \texttt{instruction\_count},
\texttt{register\_allocations}, \texttt{spills}). \texttt{run\_mips.py} acepta
\texttt{--json} para el reporte estructurado, \texttt{--max-steps N} para el
límite de instrucciones y \texttt{--delayed-branches} para activar los saltos
retardados del simulador.

% ----------------------------------------------- 15 DECISIONES Y SUPUESTOS
\section{Decisiones, supuestos y limitaciones}

\begin{itemize}[leftmargin=1.4em]
  \item \textbf{Solo se genera MIPS para programas válidos.} Léxico, sintaxis,
        semántica y TAC deben terminar sin errores; si no, el bloque
        \texttt{mips} se marca como omitido.
  \item \textbf{Valores con etiqueta de tipo.} Cada slot mide 8 bytes (carga útil
        y etiqueta). Duplica el espacio de un entero, pero permite reciclar
        temporales, imprimir cualquier valor y capturar errores de tipo
        desconocido.
  \item \textbf{Dobles por referencia.} Mantienen uniforme el tamaño del slot a
        costa de una reserva dinámica por resultado flotante calculado.
  \item \textbf{Despacho estático.} Los métodos se resuelven con el símbolo
        determinado por el análisis semántico y se llaman directamente; la
        herencia conserva campos y búsqueda de miembros, sin tabla virtual.
  \item \textbf{Closures limitados a frames activos.} Las funciones anidadas
        acceden a los frames léxicos mientras están en la pila; la gramática no
        define funciones de primera clase que escapen de su activación.
  \item \textbf{Saltos sin retardo.} Se emite un \texttt{nop} tras cada salto,
        de modo que el código es correcto con o sin la opción
        \texttt{--delayed-branches} de MARS.
  \item \textbf{Orden de evaluación.} Los argumentos se evalúan de izquierda a
        derecha y cada resultado se conserva antes de evaluar el siguiente.
\end{itemize}

\paragraph{Limitaciones conocidas.}

\begin{itemize}[leftmargin=1.4em]
  \item La asignación de registros es local: los descriptores se invalidan en
        cada etiqueta, salto y llamada, lo que produce cargas y \emph{spills}
        que una asignación global evitaría.
  \item No hay optimización: no se pliegan constantes y cada operación
        aritmética llama a una rutina auxiliar que inspecciona las etiquetas,
        aun cuando los tipos son conocidos en compilación.
  \item No hay recolector de basura: la memoria dinámica se conserva hasta que
        termina el programa.
  \item La pila de manejadores y la memoria dinámica dependen de los servicios de
        MARS; ejecutar en hardware real requeriría sustituir esos servicios.
\end{itemize}

% --------------------------------------------------- 16 RÚBRICA
\section{Correspondencia con la rúbrica}

\begin{table}[H]
\centering
\small
\begin{tabular}{@{}L{3.6cm}rL{9.4cm}@{}}
\toprule
\textbf{Componente} & \textbf{Pts} & \textbf{Evidencia en el proyecto} \\
\midrule
Llamadas y retornos & 20 &
Registros de activación físicos con enlaces de control y acceso, prólogo y
epílogo, paso de parámetros por pila, retorno en \texttt{\$v0}/\texttt{\$v1},
preservación de registros (sección~7); 8 casos de llamadas y 7 de ámbitos y
closures ejecutados en MARS (sección~13). \\
Asignación de registros & 20 &
\texttt{getReg} con descriptores, LRU, \emph{spills} y \texttt{flush} en
fronteras de control; métricas observables en el JSON y el IDE (sección~9);
casos de presión de registros y temporales vivos entre llamadas. \\
Código MIPS & 60 &
\texttt{MipsGenerator} de 1\,012 líneas que traduce todas las familias de TAC
(sección~8), entorno de ejecución, \texttt{try/catch} (sección~10), ejecución
real en MARS y paquete de código máquina (sección~11); 66 programas comparados
automáticamente. \\
\midrule
\textbf{Total} & \textbf{100} & \\
\bottomrule
\end{tabular}
\caption{Evidencia por componente evaluado.}
\end{table}

Respecto a los entregables: el repositorio conserva commits individuales por
integrante (sección~4); las baterías cubren casos exitosos, errores de
compilación, de generación y de ejecución (sección~13); la arquitectura y las
instrucciones de ejecución están documentadas (secciones~3 y 14, además de
\texttt{docs/generacion-mips.md}, \texttt{docs/arquitectura-compiscript.md} y
\texttt{README.md}); y el IDE compila, descarga y ejecuta los programas
(sección~12).

% ---------------------------------------------------- 17 CONCLUSIONES
\section{Conclusiones}

\begin{enumerate}[leftmargin=1.6em]
  \item Trabajar sobre cuádruplos estructurados, y no sobre su texto, permitió
        que el backend reutilizara toda la información de las fases anteriores:
        bastó con agregar los tipos de los operandos al TAC para seleccionar
        instrucciones y coerciones sin volver a analizar el programa.
  \item Representar cada valor con una etiqueta de tipo resolvió de una sola vez
        el reciclaje de temporales, la impresión polimórfica y la variable del
        \texttt{catch}, a cambio de duplicar el tamaño de los slots.
  \item Separar padre léxico y llamante en el enlace de acceso fue la decisión
        clave para que funciones anidadas, hermanas y recursivas encontraran su
        entorno correcto.
  \item Un \texttt{getReg} local con respaldo en pila es suficiente para producir
        código correcto en llamadas y en puntos de convergencia, y sus métricas
        hacen visibles los \emph{spills} que una asignación global eliminaría.
  \item Ejecutar el código en MARS y comparar la salida con resultados calculados
        desde el programa fuente convierte la verificación estructural de la fase
        anterior en una verificación de comportamiento: el compilador produce
        programas que realmente funcionan.
\end{enumerate}

% ---------------------------------------------------- 18 REFERENCIAS
\section{Referencias}

\begin{itemize}[leftmargin=1.4em]
  \item Aho, A., Lam, M., Sethi, R. y Ullman, J. \emph{Compilers: Principles,
        Techniques, and Tools}. 2.\textordfeminine{} edición, 2006. Secciones
        7.2.2--7.2.3 (registros de activación y secuencias de llamada),
        7.3.5--7.3.6 (enlaces de acceso) y 8.6.1--8.6.3 (descriptores de
        registros, generación de código y \texttt{getReg}).
  \item Vollmar, K. y Sanderson, P. \emph{MARS: MIPS Assembler and Runtime
        Simulator}, versión 4.5. Línea de comandos:
        \url{https://dpetersanderson.github.io/Help/MarsHelpCommand.html}.
        Servicios del sistema:
        \url{https://dpetersanderson.github.io/Help/SyscallHelp.html}.
  \item Parr, T. \emph{The Definitive ANTLR 4 Reference}. Documentación del
        \emph{target} de C++: \url{https://github.com/antlr/antlr4}.
  \item Repositorio del proyecto, rama \texttt{Compiladores2-Proyecto3-MIPS}:
        \url{https://github.com/G1LB3T0/Proyecto1_Compis_Generador_Analizador_Lexico/tree/Compiladores2-Proyecto3-MIPS}.
  \item Documentación interna del repositorio, en la carpeta \texttt{docs/}:
        \texttt{generacion-\allowbreak mips.md},
        \texttt{codigo-\allowbreak intermedio-\allowbreak tac.md},
        \texttt{arquitectura-\allowbreak compiscript.md} y
        \texttt{decisiones-\allowbreak compiscript.md}.
\end{itemize}

\end{document}
